% !TeX root = ../教材题库_全部条目_转写工作稿.tex
% Source: XJ-B1-C05-SU-P205-Y1036
% 第5章 / 小结与复习 小结与复习 / 教材 p.205 / PDF 第212页
% OCR未校对：公式与文字均可能有误；未完成数学TeX转写；禁止直接发给学生
% 原题及相邻原解截图（核对用）：../题目原图/XJ-B1-C05-SU-P205-Y1036-01.webp
% 本文件只在转写完成后作为题干使用；OCR 条目尚不是可靠题干。
\begin{bankitem}{第5章 小结与复习 第3题}{第5章 / 小结与复习 小结与复习 / 教材 p.205 / PDF 第212页}{XJ-B1-C05-SU-P205-Y1036}{OCR未校对：公式与文字均可能有误；未完成数学TeX转写；禁止直接发给学生}
\begin{ocrdraft}
[教材p.205 / PDF第212页]
\par
3．单位圆在三角函数的研究中有着重要的作用．为什么在单位圆中作出的有
\par
向线段能表示三角函数？尝试借助有向线段画出正弦、余弦、正切的三角函数图
\par
象．你能借助单位圆的对称性，利用定义推导出诱导公式吗？
\end{ocrdraft}
\end{bankitem}
